% ============================================================
%  بخش ۲ — مرور ادبیات
% ============================================================

\subsection{دیدگاه کلی و راهبرد مرور}
\label{subsec:lr_overview}

این بخش نسب فکری پژوهش حاضر را در چهار حوزه‌ی دانشی مجاور
ترسیم می\zwnj{}کند: مبانی مدیریت پروژه~(\lr{PM})،
مدیریت و انتخاب سبد پروژه~(\lr{PPM})،
بودجه\zwnj{}بندی \lr{PPM} و مدل\zwnj{}سازی جریان نقدی،
و یادگیری تقویتی~(\lr{RL}) کاربردشده برای مسائل تخصیص منابع ترتیبی.
این جریان\zwnj{}ها با هم هم پایه‌ی نظری و هم شکافی را تعریف می\zwnj{}کنند
که چارچوب پیشنهادی را انگیزه می\zwnj{}دهد.
مرور از دستورالعمل\zwnj{}های گزارش\zwnj{}دهی \lr{PRISMA~2020}
\citep{Page2021} پیروی می\zwnj{}کند تا شفافیت و بازتولیدپذیری
در شناسایی منابع، غربالگری، و تصمیمات گنجاندن را تضمین کند.

% - - - - - - - - - - - - - - - - - - - - - - - - - - - - - -
\subsubsection{پروتکل جستجو و غربالگری}
\label{subsubsec:prisma}
% - - - - - - - - - - - - - - - - - - - - - - - - - - - - - -

منابع از \lr{Google Scholar}، \lr{Semantic Scholar} و \lr{Scopus}
با استفاده از یک راهبرد جستجوی تدریجاً محدودکننده بازیابی شدند.
جستجوها از گسترده\zwnj{}ترین حوزه‌ی مرتبط (\emph{مدیریت پروژه})
آغاز شدند و به\zwnj{}ترتیب از طریق
\emph{مدیریت سبد پروژه}،
\emph{بودجه\zwnj{}بندی سبد و بهینه\zwnj{}سازی جریان نقدی}،
\emph{مدیریت ارزش کسب\zwnj{}شده در شرایط عدم\zwnj{}قطعیت}،
و \emph{یادگیری تقویتی برای زمان\zwnj{}بندی پروژه و تخصیص منابع}
پالایش شدند.
پیمایش شبکه‌ی استناد---برف\zwnj{}گلوله\zwnj{}سازی رو به عقب و رو به جلو
از مقالات لنگر با استناد بالا---برای بازیابی آثار بنیادین
که با جستجوی کلیدواژه پدیدار نشده بودند به کار رفت.

رکوردها در سه مرحله‌ی خودکار متوالی منطبق با نمودار جریان
\lr{PRISMA~2020} (پیوست~\ref{app:prisma}) غربال شدند.
خط لوله‌ی غربالگری، به\zwnj{}همراه کورپوس بایگانی\zwnj{}شده
و مصنوعات خروجی آن، به\zwnj{}عنوان یک زنجیره‌ی پردازشی بازتولیدپذیر
پیاده\zwnj{}سازی شده است که موقعیت مخزن و نگاشت مفهومی\zwnj{}سطح آن
در پیوست~\ref{app:repo_map} مستند شده است.

\medskip
\textbf{مرحله‌ی ۱ --- شناسایی و ماشین\zwnj{}خوانایی.}\;
فایل\zwnj{}های \lr{PDF} کاندیدا تحت تریاژ بیبلیوگرافیک و خوانایی
خودکار قرار گرفتند.
هر سند از نظر بازده استخراج متن (حداقل تراکم کاراکتر خوانا در صفحات)،
کیفیت واژگانی (نسبت توکن الفبایی و میانگین طول کلمه)،
و قابلیت جاسازی با تبدیل\zwnj{}گر جمله ارزیابی شد.
اسنادی که در هر معیار موفق نشدند در یک بایگانی رد اختصاصی
قرنطینه شدند، با یک هش محتوایی ثبت\zwnj{}شده تا از ورود مجدد
ارسال\zwnj{}های تکراری به خط لوله جلوگیری شود.

\medskip
\textbf{مرحله‌ی ۲ --- فیلترسازی ارتباط.}\;
رکوردهای عبورکرده از مرحله‌ی ۱ با یک امتیاز ارتباط مرکب رتبه\zwnj{}بندی شدند.
یک مؤلفه‌ی شباهت معنایی خلاصه، مقدمه و نتیجه‌گیری هر مقاله را
با یک مدل تبدیل\zwnj{}گر جمله جاسازی کرد و
حداکثر شباهت کسینوسی را در برابر دوازده بردار کوئری مشتق\zwnj{}شده از پایان\zwnj{}نامه
که حوزه\zwnj{}های دانشی هدف را پوشش می\zwnj{}دادند محاسبه کرد.
یک مؤلفه‌ی امتیازدهی کلیدواژه شمارش\zwnj{}های وزن\zwnj{}دار اصطلاحات خاص حوزه
را تخصیص داد.
دو مؤلفه با وزن مساوی ترکیب شدند تا یک امتیاز ارتباط ترکیبی
در $[0,1]$ ایجاد کنند.
مقالات به \emph{بسیار مرتبط} (امتیاز $\geq 0.45$)،
\emph{ارزشمند} (امتیاز $\geq 0.30$)،
و \emph{نامرتبط} (امتیاز $< 0.30$) تریاژ شدند،
با دو گروه اول برای استخراج ساختاریافته نگه\zwnj{}داری\zwnj{}شده.

\medskip
\textbf{مرحله‌ی ۳ --- استخراج ساختاریافته.}\;
هر مقاله‌ی نگه\zwnj{}داری\zwnj{}شده تحت استخراج دانش ساختاریافته قرار گرفت.
خلاصه، مقدمه و نتیجه‌گیری هر \lr{PDF} تجزیه شدند
و با یک پرامپت طرح\zwnj{}واره‌ی ثابت به یک مدل زبانی بزرگ ارسال شدند
که یک رکورد به\zwnj{}ازای هر مقاله تولید کرد شامل:
قاب\zwnj{}بندی مسئله، سطح سبد، نوع تصمیم،
پرداخت جریان نقدی تصادفی، مکانیزم\zwnj{}های قراردادی،
مقررات تسویه\zwnj{}ی فسخ، پارادایم مدل\zwnj{}سازی،
پارامترهای کالیبراسیون، و پرچم\zwnj{}های ارتباط\zwnj{}سنجی پایان\zwnj{}نامه.
کورپوس ساختاریافته‌ی حاصل یک داشبورد سنتز خودکار و
یک مولد جدول را هدایت کرد که ماتریس تحلیل شکاف،
نمودارهای توزیع حوزه، و نقشه‌ی ادبیات بازتولیدشده در
جدول~\ref{tab:litmap} را تولید کرد.

\medskip
خط لوله‌ی کامل---اسکریپت\zwnj{}های پردازشی، کورپوس رکورد ساختاریافته،
خروجی\zwnj{}های داشبورد، و جداول سنتز---در مخزن پروژه بایگانی شده
و یک کورپوس مرور ادبیات بازتولیدپذیر و قابل حسابرسی ماشینی
را تشکیل می\zwnj{}دهد؛ سازماندهی سطح\zwnj{}دایرکتوری آن در
پیوست~\ref{app:repo_map} شرح داده شده است.
نمودار جریان \lr{PRISMA} که رکوردها را در هر مرحله کمّی می\zwnj{}کند
در پیوست~\ref{app:prisma} ارائه شده است.

\begin{figure}[htbp]
    \centering
    % figures/screening_pipeline.tex
% TikZ only — no figure environment, no caption, no label.
% Include in section file as:
%   % figures/screening_pipeline.tex
% TikZ only — no figure environment, no caption, no label.
% Include in section file as:
%   % figures/screening_pipeline.tex
% TikZ only — no figure environment, no caption, no label.
% Include in section file as:
%   \input{figures/screening_pipeline.tex}
\begin{tikzpicture}[
    box/.style      = {rectangle, rounded corners=6pt, thick,
                       text width=7.0cm, align=center,
                       minimum height=0.85cm, inner sep=7pt,
                       font=\small},
    decision/.style = {diamond, thick, aspect=2.5,
                       minimum width=4.5cm, minimum height=0.7cm,
                       align=center, font=\scriptsize, inner sep=2pt},
    rej/.style      = {rectangle, rounded corners=6pt, thick,
                       text width=3.2cm, align=center,
                       minimum height=0.75cm, inner sep=5pt,
                       font=\scriptsize},
  ]

  %% --- Row 1 : left -> right ---
  \node[box, flowColorA] (s1)
    {\textbf{Stage 1 --- Identification \& Readability}\\
     Extraction yield $\cdot$ Lexical quality $\cdot$ Embeddability};

  \node[decision, flowColorGray, right=1.2cm of s1] (d1)
    {All three\\tests pass?};

  %% --- Row 2 : right -> left ---
  \node[box, flowColorB, below=2.4cm of d1] (s2)
    {\textbf{Stage 2 --- Relevance Filtering}\\
     Semantic similarity $+$ Keyword score $\to r \in [0,1]$};

  \node[decision, flowColorGray, left=1.2cm of s2] (d2)
    {$r \geq 0.30$?};

  %% --- Row 3 : left -> right ---
  \node[box, flowColorC, below=1.7cm of d2] (s3)
    {\textbf{Stage 3 --- Structured Extraction}\\
     JSON extraction: metadata $\cdot$ calibration $\cdot$ thesis flags};

  \node[box, flowColorD, right=1.2cm of s3] (out)
    {\textbf{Synthesis}\\
     Gap matrix $\cdot$ Domain charts $\cdot$ Literature map};

  %% --- Rejection nodes ---
  \node[rej, flowColorReject, above=0.55cm of d1] (r1)
    {Rejected\\(failed readability)};

  \node[rej, flowColorReject, above=0.55cm of d2] (r2)
    {Rejected\\(irrelevant)};

  %% --- Main snake flow ---
  \draw[arrowMain] (s1) -- (d1);
  \draw[arrowMain] (d1) -- node[right, font=\scriptsize]{Yes} (s2);
  \draw[arrowMain] (s2) -- (d2);
  \draw[arrowMain] (d2) -- node[left,  font=\scriptsize]{Yes} (s3);
  \draw[arrowMain] (s3) -- (out);

  %% --- Rejection branches ---
  \draw[arrowReject] (d1) -- node[right, font=\scriptsize]{No} (r1);
  \draw[arrowReject] (d2) -- node[left,  font=\scriptsize]{No} (r2);

\end{tikzpicture}
\begin{tikzpicture}[
    box/.style      = {rectangle, rounded corners=6pt, thick,
                       text width=7.0cm, align=center,
                       minimum height=0.85cm, inner sep=7pt,
                       font=\small},
    decision/.style = {diamond, thick, aspect=2.5,
                       minimum width=4.5cm, minimum height=0.7cm,
                       align=center, font=\scriptsize, inner sep=2pt},
    rej/.style      = {rectangle, rounded corners=6pt, thick,
                       text width=3.2cm, align=center,
                       minimum height=0.75cm, inner sep=5pt,
                       font=\scriptsize},
  ]

  %% --- Row 1 : left -> right ---
  \node[box, flowColorA] (s1)
    {\textbf{Stage 1 --- Identification \& Readability}\\
     Extraction yield $\cdot$ Lexical quality $\cdot$ Embeddability};

  \node[decision, flowColorGray, right=1.2cm of s1] (d1)
    {All three\\tests pass?};

  %% --- Row 2 : right -> left ---
  \node[box, flowColorB, below=2.4cm of d1] (s2)
    {\textbf{Stage 2 --- Relevance Filtering}\\
     Semantic similarity $+$ Keyword score $\to r \in [0,1]$};

  \node[decision, flowColorGray, left=1.2cm of s2] (d2)
    {$r \geq 0.30$?};

  %% --- Row 3 : left -> right ---
  \node[box, flowColorC, below=1.7cm of d2] (s3)
    {\textbf{Stage 3 --- Structured Extraction}\\
     JSON extraction: metadata $\cdot$ calibration $\cdot$ thesis flags};

  \node[box, flowColorD, right=1.2cm of s3] (out)
    {\textbf{Synthesis}\\
     Gap matrix $\cdot$ Domain charts $\cdot$ Literature map};

  %% --- Rejection nodes ---
  \node[rej, flowColorReject, above=0.55cm of d1] (r1)
    {Rejected\\(failed readability)};

  \node[rej, flowColorReject, above=0.55cm of d2] (r2)
    {Rejected\\(irrelevant)};

  %% --- Main snake flow ---
  \draw[arrowMain] (s1) -- (d1);
  \draw[arrowMain] (d1) -- node[right, font=\scriptsize]{Yes} (s2);
  \draw[arrowMain] (s2) -- (d2);
  \draw[arrowMain] (d2) -- node[left,  font=\scriptsize]{Yes} (s3);
  \draw[arrowMain] (s3) -- (out);

  %% --- Rejection branches ---
  \draw[arrowReject] (d1) -- node[right, font=\scriptsize]{No} (r1);
  \draw[arrowReject] (d2) -- node[left,  font=\scriptsize]{No} (r2);

\end{tikzpicture}
\begin{tikzpicture}[
    box/.style      = {rectangle, rounded corners=6pt, thick,
                       text width=7.0cm, align=center,
                       minimum height=0.85cm, inner sep=7pt,
                       font=\small},
    decision/.style = {diamond, thick, aspect=2.5,
                       minimum width=4.5cm, minimum height=0.7cm,
                       align=center, font=\scriptsize, inner sep=2pt},
    rej/.style      = {rectangle, rounded corners=6pt, thick,
                       text width=3.2cm, align=center,
                       minimum height=0.75cm, inner sep=5pt,
                       font=\scriptsize},
  ]

  %% --- Row 1 : left -> right ---
  \node[box, flowColorA] (s1)
    {\textbf{Stage 1 --- Identification \& Readability}\\
     Extraction yield $\cdot$ Lexical quality $\cdot$ Embeddability};

  \node[decision, flowColorGray, right=1.2cm of s1] (d1)
    {All three\\tests pass?};

  %% --- Row 2 : right -> left ---
  \node[box, flowColorB, below=2.4cm of d1] (s2)
    {\textbf{Stage 2 --- Relevance Filtering}\\
     Semantic similarity $+$ Keyword score $\to r \in [0,1]$};

  \node[decision, flowColorGray, left=1.2cm of s2] (d2)
    {$r \geq 0.30$?};

  %% --- Row 3 : left -> right ---
  \node[box, flowColorC, below=1.7cm of d2] (s3)
    {\textbf{Stage 3 --- Structured Extraction}\\
     JSON extraction: metadata $\cdot$ calibration $\cdot$ thesis flags};

  \node[box, flowColorD, right=1.2cm of s3] (out)
    {\textbf{Synthesis}\\
     Gap matrix $\cdot$ Domain charts $\cdot$ Literature map};

  %% --- Rejection nodes ---
  \node[rej, flowColorReject, above=0.55cm of d1] (r1)
    {Rejected\\(failed readability)};

  \node[rej, flowColorReject, above=0.55cm of d2] (r2)
    {Rejected\\(irrelevant)};

  %% --- Main snake flow ---
  \draw[arrowMain] (s1) -- (d1);
  \draw[arrowMain] (d1) -- node[right, font=\scriptsize]{Yes} (s2);
  \draw[arrowMain] (s2) -- (d2);
  \draw[arrowMain] (d2) -- node[left,  font=\scriptsize]{Yes} (s3);
  \draw[arrowMain] (s3) -- (out);

  %% --- Rejection branches ---
  \draw[arrowReject] (d1) -- node[right, font=\scriptsize]{No} (r1);
  \draw[arrowReject] (d2) -- node[left,  font=\scriptsize]{No} (r2);

\end{tikzpicture}
    \caption{خط لوله‌ی غربالگری سه\zwnj{}مرحله\zwnj{}ای تراز\zwnj{}یافته با \lr{PRISMA}،
      به\zwnj{}صورت جریان مارپیچی ترتیب\zwnj{}یافته.
      جعبه\zwnj{}ها مراحل پردازش خودکار را نشان می\zwnj{}دهند؛
      الماس\zwnj{}ها دروازه\zwnj{}های قبول/رد را نشان می\zwnj{}دهند؛
      فلش\zwnj{}های نقطه\zwnj{}چین مسیریابی رد را نشان می\zwnj{}دهند.
      جزئیات پیاده\zwnj{}سازی و موقعیت\zwnj{}های مخزن در
      پیوست~\ref{app:repo_map} ارائه شده است.}
    \label{fig:screening_pipeline}
\end{figure}

% ------------------------------------------------------------
\subsection{مدیریت پروژه: مفاهیم بنیادی}
\label{subsec:lr_pm}
% ------------------------------------------------------------

مطالعه‌ی آکادمیک مدیریت پروژه بر کنترل زمان\zwnj{}بندی
(\lr{CPM}/\lr{PERT})، مدیریت ارزش کسب\zwnj{}شده~(\lr{EVM})،
و مدل\zwnj{}سازی ریسک به\zwnj{}عنوان ستون\zwnj{}های کمّی اصلی خود استوار است
\citep{Fleming2010, PMI2021}.
\lr{EVM} شاخص عملکرد هزینه و شاخص عملکرد زمان\zwnj{}بندی را
به\zwnj{}عنوان سیگنال\zwnj{}های تشخیصی بلادرنگ معرفی کرد که
هزینه‌ی بودجه\zwnj{}شده‌ی کار انجام\zwnj{}شده را با مخارج واقعی
و زمان\zwnj{}بندی برنامه\zwnj{}ریزی\zwnj{}شده مرتبط می\zwnj{}سازند---
واژگان عملکردی که مدل حاضر مستقیماً از آن\zwnj{}ها به ارث می\zwnj{}برد.
مدل حاضر صورت\zwnj{}بندی زمان\zwnj{}بندی کسب\zwnj{}شده
از شاخص عملکرد زمان\zwnj{}بندی را می\zwnj{}پذیرد \citep{Lipke2003}،
که پس از تاریخ پایان برنامه\zwnj{}ریزی\zwnj{}شده نیز به تنزل ادامه می\zwnj{}دهد
به\zwnj{}جای اینکه هنگامی که ارزش برنامه\zwnj{}ریزی\zwnj{}شده به بودجه‌ی کامل می\zwnj{}رسد
منجمد شود---یک محدودیت شناخته\zwnj{}شده‌ی نسبت کلاسیک
در دوره\zwnj{}های بازیابی پس از پایان برنامه\zwnj{}ریزی\zwnj{}شده
\citep{Lipke2003, Henderson2004}.
شاخص عملکرد تا تکمیل به\zwnj{}عنوان مکمل آینده\zwnj{}نگر گنجانده شده
و کارایی هزینه‌ی موردنیاز در کار باقی\zwnj{}مانده
برای پایان\zwnj{}یافتن در بودجه را اندازه\zwnj{}گیری می\zwnj{}کند.
برآورد در تکمیل از یک مخرج مرکب هزینه\zwnj{}و\zwnj{}زمان\zwnj{}بندی استفاده می\zwnj{}کند؛
\citet{christensen1993cost} در ۱۵۵ قرارداد نشان داد که این صورت\zwnj{}بندی
دقیق\zwnj{}ترین و پایدارترین پیش\zwnj{}بینی\zwnj{}های هزینه\zwnj{}در\zwnj{}تکمیل را تولید می\zwnj{}کند،
با افزایش اتکاپذیری فراتر از نقطه‌ی ۲۰\% تکمیل \citep{Lipke2003}.
توسعه\zwnj{}های تصادفی \lr{EVM} \citep{Acebes2015, Batselier2017}
و چارچوب\zwnj{}های شبیه\zwnj{}سازی مونت\zwnj{}کارلو \citep{Kwak2007}
عدم\zwnj{}قطعیت عملکرد را در سطح پروژه مدل کرده\zwnj{}اند،
ثابت کرده\zwnj{}اند که شاخص\zwnj{}های عملکرد از توزیع\zwnj{}های
مشخص\zwnj{}سازی\zwnj{}پذیر تجربی در طول عمر پروژه پیروی می\zwnj{}کنند
و پایه‌ی کالیبراسیون پارامتر بازده تصادفی در محیط حاضر را فراهم می\zwnj{}کنند.

تحلیل منحنی\lr{-S}، که توسط \citet{Kenley2009} برای جریان\zwnj{}های
نقدی ساخت\zwnj{}وساز رسمی\zwnj{}سازی شده،
نشان داد که مصرف بودجه از پروفیل\zwnj{}های سیگموئیدی مشخص پیروی می\zwnj{}کند
که پارامترهای شکل آن\zwnj{}ها به نوع پروژه و فرم قرارداد وابسته است.
تابع بتا ناقص نرمال\zwnj{}شده یک پارامتری\zwnj{}سازی خانواده‌ی منفرد انعطاف\zwnj{}پذیر
از این پروفیل فراهم می\zwnj{}کند \citep{barraza2004probabilistic}
و مستقیماً در مدل مخارج برنامه\zwnj{}ریزی\zwnj{}شده‌ی صورت\zwnj{}بندی حاضر اتخاذ شده است.

شواهد تجربی در مقیاس بزرگ درباره‌ی عملکرد پروژه توسط
\citet{Merrow2011} از طریق پایگاه داده‌ی تحلیل مستقل پروژه،
و توسط \citet{Flyvbjerg2003, Flyvbjerg2018} از طریق
مجموعه\zwnj{}داده\zwnj{}های تجاوز از هزینه‌ی زیرساخت در دهه\zwnj{}ها و جغرافیاهای مختلف
ارائه شده است.
این مجموعه\zwnj{}داده\zwnj{}ها هم کالیبراسیون توزیع\zwnj{}های عدم\zwnj{}قطعیت
و هم معیارهایی که پارامترهای محیط مصنوعی در برابر آن\zwnj{}ها اعتبارسنجی می\zwnj{}شوند
را آگاه می\zwnj{}سازند.
معیارهای تجربی تکمیلی از پایگاه\zwnj{}های داده‌ی مهندسی هزینه‌ی
\lr{AACE International} و \lr{Construction Industry Institute}
استخراج شده\zwnj{}اند \citep{christensen1993cost}.

ساختارهای حاکمیت قراردادی---به\zwnj{}ویژه فرم\zwnj{}های \lr{FIDIC}
(کتاب\zwnj{}های قرمز، زرد و نقره\zwnj{}ای)---
تعامل مالی بین کارفرما و پیمانکار را از طریق
برنامه\zwnj{}های پرداخت نقطه‌ی عطف، مکانیزم\zwnj{}های پیش\zwnj{}پرداخت و بازیابی،
و بندهای حسن انجام کار تجویز می\zwnj{}کنند \citep{FIDIC2017}.
تحت بند فرعی ۱۴.۲ \lr{FIDIC}، پیش\zwnj{}پرداخت یک وام بسیج
بدون بهره است که از طریق کسورات درصدی از گواهی\zwnj{}های پرداخت
بین\zwnj{}دوره\zwnj{}ای بازیابی می\zwnj{}شود؛
آستانه‌ی شروع و نرخ بازیابی به\zwnj{}عنوان شرایط خاص
پروژه\zwnj{}محور مشخص می\zwnj{}شوند \citep{FIDIC2017}.
مدل حاضر این مکانیزم را با نرخ\zwnj{}های پیش\zwnj{}پرداخت و بازیابی
خاص هر پروژه پارامتری\zwnj{}سازی می\zwnj{}کند و
کل طیف شرایط خاص موجود در عمل را در بر می\zwnj{}گیرد.

% ------------------------------------------------------------
\subsection{مدیریت سبد پروژه: انتخاب و حاکمیت}
\label{subsec:lr_ppm}
% ------------------------------------------------------------

مدیریت سبد پروژه مسئله‌ی تخصیص منابع محدود یک سازمان
در میان مجموعه\zwnj{}ای رقیب از پروژه\zwnj{}های کاندیدا
برای بیشینه\zwnj{}سازی ارزش راهبردی را رسمی می\zwnj{}کند
\citep{Archer1999, Cooper2001}.
ادبیات دو لایه‌ی تصمیم\zwnj{}گیری تودرتو را متمایز می\zwnj{}کند:
\emph{انتخاب سبد}---بهینه\zwnj{}سازی باینری یا عدد\zwnj{}صحیح مختلط
از اینکه کدام پروژه\zwnj{}ها گنجانده شوند---
و \emph{حاکمیت سبد}---نظارت عملیاتی مستمر بر پروژه\zwnj{}های تأییدشده و فعال.

مدل\zwnj{}های انتخاب سبد به\zwnj{}طور گسترده توسعه یافته\zwnj{}اند.
چارچوب\zwnj{}های میانگین\zwnj{}واریانس به سبک مارکویتز
اقتباس\zwnj{}شده از نظریه‌ی سبد مالی \citep{Markowitz1952}
برای بدل\zwnj{}کردن \lr{NPV} مورد انتظار در برابر واریانس
بر سبدهای پروژه اعمال شده\zwnj{}اند \citep{Eilat2006}.
توسعه\zwnj{}های تحلیل تصمیم\zwnj{}گیری چندمعیاره، همسویی راهبردی،
رتبه\zwnj{}بندی\zwnj{}های ریسک، و وابستگی\zwnj{}های منابع را در بر می\zwnj{}گیرند
\citep{Martinsuo2007, Killen2008}.
رویکردهای فراابتکاری نمونه\zwnj{}های ترکیباتی بزرگ را
پوشش می\zwnj{}دهند \citep{Zhang2020}.
انتخاب سبد مقیدشده به منابع تحت سهمیه\zwnj{}بندی سرمایه توسط
\citet{kavadias2003resource} مطالعه شده است
که قید عدم\zwnj{}استقراض او مستقیماً در صورت\zwnj{}بندی حاضر اتخاذ شده است.

پژوهش حاضر صریحاً خروجی انتخاب سبد را به\zwnj{}عنوان ورودی برون\zwnj{}زا می\zwnj{}پذیرد.
مسئله‌ی انتخاب به\zwnj{}عنوان حل\zwnj{}شده در بالادست تلقی می\zwnj{}شود؛
تمرکز کاملاً بر لایه‌ی حاکمیت عملیاتی است که
یک سبد متعهدشده دریافت می\zwnj{}کند و باید آن را در طول اجرا مدیریت کند.
این تجزیه منطبق با مدل حاکمیت دومرحله\zwnj{}ای \citet{Turner2014}
و استاندارد \lr{PMI} برای مدیریت سبد \citep{PMI2017} است
که انتخاب راهبردی را از نظارت اجرای عملیاتی جدا می\zwnj{}کند.
با این کار، این پژوهش مستقیماً بر ادبیات انتخاب بنا می\zwnj{}شود
و در عین حال مرز پیشین را به حوزه‌ی حاکمیت پیش می\zwnj{}راند
که توجه تحلیلی رسمی به مراتب کمتری دریافت کرده است.

% ------------------------------------------------------------
\subsection{بودجه\zwnj{}بندی \lr{PPM}، مدل\zwnj{}سازی جریان نقدی و عدم\zwnj{}قطعیت}
\label{subsec:lr_cashflow}
% ------------------------------------------------------------

در حاکمیت سبد، تخصیص بودجه میان پروژه\zwnj{}های فعال در طول زمان
یک تصمیم عملیاتی پرتکرار و پرخطر را نمایندگی می\zwnj{}کند.
با وجود اهمیت عملی آن، این مسئله‌ی خاص
توجه رسمی محدودتری نسبت به انتخاب سبد دریافت کرده است.
آثار موجود مدل\zwnj{}سازی جریان نقدی سطح پروژه
\citep{Kenley2009, Park2004}،
پیش\zwnj{}بینی هزینه\zwnj{}در\zwnj{}تکمیل پروژه‌ی منفرد مبتنی بر \lr{EVM}
\citep{Kim2002, Acebes2015}،
و تراز\zwnj{}سازی منابع چندپروژه\zwnj{}ای \citep{Lova2000, Peteghem2010}
را پوشش می\zwnj{}دهند.
این سهم\zwnj{}ها بودجه\zwnj{}ها را به\zwnj{}عنوان قیود قطعی تلقی می\zwnj{}کنند
نه به\zwnj{}عنوان متغیرهای تصمیم پویا تابع
بهینه\zwnj{}سازی ترتیبی سطح سبد.

عدم\zwnj{}قطعیت جریان نقدی در مدل حاضر از یک منبع واحد ناشی می\zwnj{}شود:
بازده بودجه\zwnj{}به\zwnj{}پیشرفت، که حاکم است بر اینکه
یک تخصیص بودجه‌ی معین در هر دوره چه مقدار
پیشرفت فیزیکی ایجاد می\zwnj{}کند \citep{Acebes2015, Batselier2017}.
بهره\zwnj{}وری نامطمئن تعیین می\zwnj{}کند که آیا و چه زمانی هر آستانه‌ی نقطه‌ی عطف
رد می\zwnj{}شود، که به نوبه‌ی خود زمان\zwnj{}بندی پرداخت نقطه‌ی عطف قراردادی ثابت
را تعیین می\zwnj{}کند.
تأخیرهای پرداخت بین تأیید و دریافت به\zwnj{}عنوان یک منبع عدم\zwnj{}قطعیت مستقل
مدل نمی\zwnj{}شوند؛ کانال عملکردمحور منبع ساختاری اصلی عدم\zwnj{}قطعیت زمان\zwnj{}بندی
جریان نقدی در قراردادهای پرداخت نقطه‌ی عطف است
و جداسازی آن شناسایی تمیزتر مسئله‌ی تخصیص را
بدون از دست دادن قاب\zwnj{}بندی اصلی عدم\zwnj{}قطعیت ممکن می\zwnj{}سازد.
سایر منابع عدم\zwnj{}قطعیت جریان نقدی---تغییرات حوزه،
شوک\zwnj{}های بازار خارجی \citep{Doloi2012}---به\zwnj{}عنوان برون\zwnj{}زا تلقی می\zwnj{}شوند.

برنامه\zwnj{}ریزی تصادفی \citep{Birge2011} چارچوب \lr{OR} کانونی
برای بهینه\zwnj{}سازی در شرایط چنین عدم\zwnj{}قطعیتی را فراهم می\zwnj{}کند
و نوع دومرحله\zwnj{}ای آن برای زمان\zwnj{}بندی پروژه
تحت عدم\zwnj{}قطعیت منابع اعمال شده است \citep{Lova2000}.
با این حال، برنامه\zwnj{}های تصادفی چندمرحله\zwnj{}ای برای
تخصیص بودجه‌ی ترتیبی سطح سبد با انفجار درخت سناریو
با افزایش تعداد پروژه\zwnj{}ها و دوره\zwnj{}ها مواجه می\zwnj{}شوند،
همان\zwnj{}طور که در بخش~\ref{subsec:theoretical_ub} رسمی\zwnj{}سازی شده است.
این مسئله‌ی سختی صورت\zwnj{}بندی \lr{MDP} پژوهش حاضر را توجیه می\zwnj{}کند.

شکاف در ادبیات دقیق است: هیچ چارچوب تثبیت\zwnj{}شده\zwnj{}ای
تصمیم تخصیص بودجه‌ی سطح سبد را به\zwnj{}عنوان یک مسئله‌ی
بهینه\zwnj{}سازی تصادفی ترتیبی مدل نمی\zwnj{}کند که در آن
یک تصمیم\zwnj{}گیرنده‌ی مرکزی منابع را میان چند پروژه‌ی فعال
تخصیص می\zwnj{}دهد که مسیرهای ورودی نقدی آن\zwnj{}ها از طریق
بازده تصادفی بودجه\zwnj{}به\zwnj{}پیشرفت نامطمئن هستند،
تابع پنجره\zwnj{}های تأیید نقطه‌ی عطف قراردادی
که سرعت تبدیل پیشرفت به ورودی\zwnj{}های نقدی را محدود می\zwnj{}کنند.

تابع پاداش محیط حاضر جریان نقدی خالص تنزیل\zwnj{}شده‌ی ناب است؛
هیچ جریمه\zwnj{}ی انحراف در هدف وارد نمی\zwnj{}شود.
این یک انتخاب طراحی عمدی است: پیامدهای انحراف پایدار---
فسخ و از دست دادن تمام پرداخت\zwnj{}های نقطه‌ی عطف آینده
و مبالغ نگه\zwnj{}داری\zwnj{}شده---برای نامطلوب\zwnj{}کردن انحراف
بدون نیاز به یک جمله‌ی جریمه‌ی جداگانه کافی هستند،
و گنجاندن آن یک وزن اسکالارسازی دلبخواهانه
بدون کالیبراسیون طبیعی از داده\zwnj{}های قراردادی معرفی می\zwnj{}کرد \citep{Ng1999}.
ضریب تنزیل نه تنها ترجیح زمانی بلکه
هزینه‌ی فرصت سرمایه‌ی گره\zwnj{}خورده را در یک محیط رقابتی رمزگذاری می\zwnj{}کند
که در آن پروژه\zwnj{}های آینده در حالی که پروژه\zwnj{}های جاری ناتمام مانده\zwnj{}اند
از دست می\zwnj{}روند \citep{Damodaran2012}،
و یک لنگر کالیبراسیون مشخص از نرخ\zwnj{}های استقرار مجدد خط لوله‌ی مشاهده\zwnj{}شده‌ی شرکت
یا معیارهای هزینه‌ی سرمایه‌ی رایج به آن می\zwnj{}دهد.
یک پیامد این است که یک سیاست خوب\zwnj{}عملکرد ممکن است
مسیرهای تخصیصی ایجاد کند که مخارج را حول آستانه\zwnj{}های تأیید نقطه‌ی عطف
متمرکز می\zwnj{}کند، و راهبرد مالی مستند پیمانکار در قراردادهای پرداخت نقطه‌ی عطف
را بازتاب می\zwnj{}دهد \citep{Kenley2009, Kaka1993, Arditi2005}.

% ------------------------------------------------------------
\subsection{عملکرد ضعیف پیمانکار، دوره\zwnj{}های مهلت و فسخ قرارداد}
\label{subsec:lr_interventions}
% ------------------------------------------------------------

هنگامی که یک پیمانکار در رعایت آستانه\zwnj{}های عملکرد زمان\zwnj{}بندی
یا هزینه ناکام می\zwnj{}ماند، قراردادهای بین\zwnj{}المللی \lr{EPC}
یک توالی تشدید ساختاریافته قبل از اعمال فسخ را تجویز می\zwnj{}کنند.
تحت بند فرعی ۱۵.۱ \lr{FIDIC}، کارفرما باید یک اخطار رسمی نقض
صادر کند و یک دوره‌ی مهلت---معمولاً ۲۸ روز---اجازه دهد
که طی آن پیمانکار موظف است بهبود قابل اندازه\zwnj{}گیری را نشان دهد
پیش از اینکه کارفرما بتواند اقدام به فسخ کند
\citep{FIDIC2017, Ndekugri2008}.
این دوره‌ی مهلت یک مهلت اختیاری نیست؛
یک تعهد قراردادی است که رعایت\zwnj{}نکردن آن
فسخ نادرست را تشکیل می\zwnj{}دهد.
مطالعات تجربی اختلافات بین\zwnj{}المللی \lr{EPC} تأیید می\zwnj{}کنند
که عدم رعایت دوره‌ی مهلت از پرتکرارترین مبانی مورد اعتراض
در ادعاهای فسخ پیمانکار است \citep{Ndekugri2008, Ennis2021}.

ادبیات مدیریت پروژه‌ی رفتاری شرایطی را مطالعه کرده است
که تحت آن\zwnj{}ها سازمان\zwnj{}ها در خاتمه دادن به پروژه\zwnj{}های با عملکرد ضعیف
در نقطه‌ی اقتصادی عقلانی ناکام می\zwnj{}مانند.
تشدید تعهد \citep{Staw1987}---تمایل به ادامه‌ی سرمایه\zwnj{}گذاری
در پروژه\zwnj{}های در حال شکست فراتر از آنچه ارزش مورد انتظار توجیه می\zwnj{}کند---
در برنامه\zwnj{}های ساخت\zwnj{}وساز، فناوری اطلاعات و زیرساخت مستند شده است.
نظریه‌ی اختیارات واقعی \citep{Dixit1994, Huchzermeier2001}
مکمل هنجاری را فراهم می\zwnj{}کند: فسخ اعمال یک اختیار رهاسازی است
که ارزش آن از \lr{NPV} ادامه هنگامی که عملکرد
از یک آستانه‌ی قابل بازیابی عبور کرده است بیشتر می\zwnj{}شود.
این جریان\zwnj{}ها با هم ثابت می\zwnj{}کنند که تصمیم برای ادامه یا فسخ
یک پروژه‌ی با عملکرد ضعیف
نه صرفاً قراردادی است نه صرفاً مالی---
بلکه شامل سیگنال\zwnj{}های عملکرد، تعهدات قراردادی،
و زمان\zwnj{}بندی همزمان می\zwnj{}شود.

مدل حاضر هر سه بعد را در یک مکانیزم محیط واحد
تحت حاکمیت دو شرط ادغام می\zwnj{}کند.
یک نقض تأخیر پیشرفت زمانی آتش می\zwnj{}گیرد که تأخیر بین
پیشرفت برنامه\zwnj{}ریزی\zwnj{}شده و واقعی از تحمل قراردادی فراتر رود،
مبتنی بر اختیار مهندس \lr{FIDIC} برای صدور اخطار هنگامی که
پیشرفت به\zwnj{}طور قابل توجهی از برنامه عقب می\zwnj{}افتد \citep{FIDIC2017}.
یک نقض مهلت سخت به\zwnj{}صورت غیرمشروط در تجاوز از پایان
مجاز قراردادی آتش می\zwnj{}گیرد و فسخ فوری را
صرف\zwnj{}نظر از وضعیت شمارنده‌ی دوره‌ی مهلت فعال می\zwnj{}کند.
وقتی یک نقض تأخیر پیشرفت فعال است،
شمارنده‌ی دوره‌ی مهلت به اندازه‌ی یک کاهش می\zwnj{}یابد؛
تنها زمانی که هیچ نقضی وجود ندارد به مقدار اولیه‌ی خود بازنشینی می\zwnj{}کند.
فسخ زمانی اجرا می\zwnj{}شود که شمارنده به صفر برسد.
این مکانیزم همزمان در حقوق قراردادی \lr{FIDIC}
\citep{FIDIC2017, Ndekugri2008}،
قاب\zwnj{}بندی اختیارات واقعی از آستانه‌ی رهاسازی
\citep{Dixit1994, Huchzermeier2001}،
اعتبارسنجی ادبیات \lr{EVM} از شاخص\zwnj{}های عملکرد به\zwnj{}عنوان سیگنال\zwnj{}های
پیش\zwnj{}بینی\zwnj{}پایدار \citep{christensen1993est, Lipke2003}،
و شواهد ادبیات تشدید تعهد مبنی بر اینکه سیگنال\zwnj{}های عملکرد
باید پایدار---نه گذرا---باشند قبل از اینکه فسخ توجیه شود
\citep{Staw1987} ریشه دارد.

شمارنده‌ی دوره‌ی مهلت بخشی از حالت قابل مشاهده است
و به سیاست یک سیگنال پیش\zwnj{}بینانه می\zwnj{}دهد
که واقعیت قراردادی را بازتاب می\zwnj{}دهد
بدون اینکه کنترل اختیاری بر نتیجه‌ی فسخ را به آن بدهد.
از دیدگاه مدیر سبد، این سیگنال یک عملکرد دوگانه دارد:
پروژه\zwnj{}های در حال رسیدن به شکست قراردادی را علامت می\zwnj{}زند
و در عین حال نقطه\zwnj{}ای را که تخصیص مستمر
بازدهی کاهشی نسبت به ارزش تسویه‌ی کار
از قبل انجام\zwnj{}شده ایجاد می\zwnj{}کند مشخص می\zwnj{}کند.
یک مدیر عقلانی که شمارنده\zwnj{}ای را می\zwnj{}بیند که در یک محیط
هزینه‌ی فرصت بالا به سمت صفر نزدیک می\zwnj{}شود
تشخیص می\zwnj{}دهد که ارزش واقعی پرداخت\zwnj{}های نقطه‌ی عطف آینده
در حال فرسایش است در حالی که ارزش تسویه‌ی کار در دست
به لحاظ اسمی ثابت است---
یک انگیزه‌ی طبیعی تخصیص مجدد ایجاد می\zwnj{}کند که مدل
از طریق تعامل ضریب تنزیل و مکانیزم\zwnj{}های تسویه‌ی فسخ ثبت می\zwnj{}کند،
بدون نیاز به مدل\zwnj{}سازی صریح تورم
\citep{Dixit1994, FIDIC2017}.

% ------------------------------------------------------------
\subsection{یادگیری تقویتی برای تخصیص منابع ترتیبی}
\label{subsec:lr_rl}
% ------------------------------------------------------------

یادگیری تقویتی یک چارچوب اصولی برای تصمیم\zwnj{}گیری ترتیبی
در شرایط عدم\zwnj{}قطعیت فراهم می\zwnj{}کند \citep{Sutton2018}.
صورت\zwnj{}بندی \lr{MDP}---حالت، عمل، گذار، پاداش---
به\zwnj{}صورت طبیعی بر مسئله‌ی تخصیص بودجه‌ی سبد منطبق می\zwnj{}شود:
مدیر سبد به\zwnj{}عنوان یک سیاست مدل می\zwnj{}شود که حالت جاری را مشاهده،
یک تخصیص متعهد، و یک پاداش جریان نقدی خالص تنزیل\zwnj{}شده دریافت می\zwnj{}کند،
در حالی که محیط به\zwnj{}صورت تصادفی بر اساس تحقق\zwnj{}های بازده گذار می\zwnj{}کند.

\lr{RL} برای زمان\zwnj{}بندی پروژه \citep{Peng2021, Liu2023}،
مسائل زمان\zwnj{}بندی پروژه‌ی مقیدشده به منابع \citep{Tassel2023}،
و مدیریت موجودی زنجیره‌ی تأمین \citep{Giannoccaro2002} اعمال شده است.
در هر حوزه، \lr{RL} مزایایی نسبت به روش\zwnj{}های کلاسیک \lr{OR}
نشان می\zwnj{}دهد وقتی فضاهای حالت بزرگ هستند،
دینامیک گذار تصادفی است، مسئله در طول افق تصمیم\zwnj{}گیری غیرایستا است،
و طول افق در آغاز برنامه\zwnj{}ریزی نامطمئن است.
اینها دقیقاً مشخصات مسئله‌ی تخصیص بودجه‌ی سبد هستند.

\lr{API} گیمنازیوم \citep{Towers2024} رابط استاندارد \textit{de facto}
برای توسعه‌ی محیط \lr{RL} را فراهم می\zwnj{}کند
و ادغام مستقیم با کتابخانه\zwnj{}های بهینه\zwnj{}سازی سیاست مدرن را ممکن می\zwnj{}سازد.
تصادفی\zwnj{}سازی حوزه \citep{Tobin2017}---تصادفی\zwnj{}سازی پارامترهای محیط
در طول اپیزودها در حین آموزش---به\zwnj{}صورت بومی از طریق
دیکشنری پیکربندی پشتیبانی می\zwnj{}شود
و یک سیاست آموزش\zwnj{}دیده را قادر می\zwnj{}سازد بدون آموزش مجدد
بر ساختارهای قراردادی و انواع پروژه‌ی متنوع تعمیم دهد.

کاربرد مستقیم \lr{RL} برای تخصیص بودجه‌ی سطح سبد
تحت مکانیزم\zwnj{}های جریان نقدی قراردادی
از ادبیات غایب مانده است.
پژوهش حاضر اولین صورت\zwnj{}بندی \lr{MDP} از این مسئله
و اولین اثبات تجربی را ارائه می\zwnj{}دهد که اکتشافی\zwnj{}های قطعی حریصانه
به\zwnj{}صورت سیستماتیک با رشد پیچیدگی سبد از یک بهینه‌ی اطلاعات\zwnj{}کامل
ضعیف\zwnj{}تر عمل می\zwnj{}کنند و یک هدف بهبود کمّی\zwnj{}شده
برای سیاست\zwnj{}های یادگیری\zwnj{}شده‌ی آینده تعیین می\zwnj{}کند.

% ------------------------------------------------------------
\subsection{ترکیب: جایگاه\zwnj{}یابی سهم حاضر}
\label{subsec:lr_synthesis}
% ------------------------------------------------------------

جدول~\ref{tab:litmap} آثار نماینده را در چهار جریان ادبیات
در برابر ابعاد مدل\zwnj{}سازی کلیدی چارچوب حاضر نگاشت می\zwnj{}کند،
همان\zwnj{}طور که توسط خط لوله‌ی سنتز خودکار
شرح\zwnj{}داده\zwnj{}شده در بخش~\ref{subsubsec:prisma} تولید شده است.
سه شکاف ساختاری از این نگاشت پدیدار می\zwnj{}شوند.

\textbf{شکاف ۱ --- انتخاب در برابر حاکمیت عملیاتی.}\;
گذار از \emph{انتخاب} سبد به \emph{حاکمیت عملیاتی} سبد
به\zwnj{}عنوان یک مسئله‌ی تصمیم\zwnj{}گیری ترتیبی مدل نشده است.
ادبیات انتخاب مدیریت پس از انتخاب را کنار می\zwnj{}گذارد؛
ادبیات زمان\zwnj{}بندی پروژه وابستگی\zwnj{}های متقابل سبد
و قید سرمایه‌ی مشترک را نادیده می\zwnj{}گیرد.
هیچ اثر موجودی هر دو را همزمان پوشش نمی\zwnj{}دهد.

\textbf{شکاف ۲ --- انتشار عدم\zwnj{}قطعیت سطح پروژه در مقابل سطح سبد.}\;
مدل\zwnj{}های تصادفی موجود از عدم\zwnj{}قطعیت جریان نقدی پروژه
در سطح پروژه‌ی منفرد عمل می\zwnj{}کنند
و عدم\zwnj{}قطعیت را به یک معماری تصمیم\zwnj{}گیری سطح سبد
تحت حاکمیت یک قید بودجه‌ی مشترک منتقل نمی\zwnj{}کنند.
حلقه‌ی بازخورد بین مسیر بازده تصادفی یک پروژه
و حق دسترسی به تأمین مالی سطح سبد آن بدون مدل باقی مانده است.

\textbf{شکاف ۳ --- غیاب مکانیزم\zwnj{}های فسخ درون\zwnj{}زا.}\;
مدل\zwnj{}سازی مشترک یک شرط نقض تأخیر پیشرفت ردیابی\zwnj{}شده
در برابر تحمل قراردادی، یک شمارنده‌ی دوره‌ی مهلت کاهشی
که خروج از دوره‌ی مهلت را اداره می\zwnj{}کند،
یک شرط نقض مهلت سخت مستقل که فسخ فوری را
صرف\zwnj{}نظر از وضعیت شمارنده‌ی دوره‌ی مهلت فعال می\zwnj{}کند،
و یک جریان نقدی تسویه‌ی فسخ که ارزش کسب\zwnj{}شده را
در برابر پرداخت\zwnj{}های تجمعی دریافت\zwnj{}شده تطبیق می\zwnj{}دهد،
در یک چارچوب تصمیم\zwnj{}گیری ترتیبی سطح سبد
که تخصیص بودجه تنها متغیر کنترلی است جاسازی نشده است.

\medskip
این شکاف\zwnj{}ها اشتباهات منفرد نیستند؛
آن\zwnj{}ها یک پیشرفت ناقص در مسیر
یک سیستم پشتیبانی تصمیم سطح سبد قابل استقرار را منعکس می\zwnj{}کنند.
شکل~\ref{fig:maturity} این پیشرفت را
به\zwnj{}صورت سه گام متوالی جای می\zwnj{}دهد.

\textbf{گام ۱}---زمینه\zwnj{}سازی حوزه---به\zwnj{}خوبی توسط ادبیات
مرورشده‌ی بالا تثبیت شده است:
مکانیزم\zwnj{}های قراردادی، سیگنال\zwnj{}های عملکرد \lr{EVM}،
پروفیل\zwnj{}های هزینه‌ی منحنی\lr{-S}، و مدل\zwnj{}های انتخاب سبد
به\zwnj{}طور دسته\zwnj{}جمعی ورودی\zwnj{}ها و شرایط مرزی را فراهم می\zwnj{}کنند
که هر مدل حاکمیت عملیاتی به آن\zwnj{}ها نیاز دارد.

\textbf{گام ۲}---صورت\zwnj{}بندی مسئله‌ی قابل حل---شکافی است که
این مقاله می\zwnj{}بندد.
سهم یک محیط \lr{MDP} تأییدشده و پیکربندی\zwnj{}پذیر است
که چرخه‌ی کامل جریان نقدی \lr{FIDIC}
و یک مکانیزم فسخ دوشرطی را رمزگذاری می\zwnj{}کند،
به\zwnj{}صورت یک بسته‌ی سازگار با \lr{Gymnasium}
با یک دفتر کل حسابرسی \lr{SQLite} کامل منتشر شده.
اکتشافی\zwnj{}های قطعی حریصانه در برابر یک بهینه‌ی اطلاعات\zwnj{}کامل
معیار می\zwnj{}شوند و یک شکاف بهینگی کمّی\zwnj{}شده ایجاد می\zwnj{}کنند
که به\zwnj{}صورت یکنواخت با پیچیدگی سبد رشد می\zwnj{}کند
و هدف عملکردی برای عامل\zwnj{}های یادگیری\zwnj{}شده‌ی آینده را تعیین می\zwnj{}کند.

\textbf{گام ۳}---استقرار اعتبارسنجی\zwnj{}شده، شامل آموزش عامل \lr{RL}
و ادغام با سیستم\zwnj{}های اطلاعات مدیریت پروژه‌ی متداول
مانند \lr{Oracle Primavera~P6}---کار آینده باقی می\zwnj{}ماند
و مشروط به محیطی است که گام ۲ اکنون تحویل می\zwnj{}دهد.

\begin{figure}[htbp]
    \centering
    % figures/maturity.tex
% TikZ only — no figure environment, no caption, no label.
% Include in section file as:
%   % figures/maturity.tex
% TikZ only — no figure environment, no caption, no label.
% Include in section file as:
%   % figures/maturity.tex
% TikZ only — no figure environment, no caption, no label.
% Include in section file as:
%   \input{figures/maturity.tex}
\begin{tikzpicture}[
    node distance  = 0cm,
    stepbox/.style = {rectangle, draw, rounded corners=3pt,
                      minimum width=3.6cm, minimum height=2.8cm,
                      align=center, font=\small, text width=3.2cm},
    arrow/.style   = {->, >=stealth, thick},
  ]

  \node[stepbox, fill=gray!15] (s1) at (0,0)
    {\textbf{Step 1}\\[4pt]
     Domain Grounding\\[6pt]
     \scriptsize
     Contract mechanics\\
     EVM signals\\
     S-curve profiles\\
     Portfolio selection};

  \node[stepbox, fill=gray!8, right=0.6cm of s1] (s2)
    {\textbf{Step 2}\\[1pt]
     \emph{(this work)}\\[4pt]
     Problem Formulation\\[6pt]
     \scriptsize
     MDP environment\\
     FIDIC cash-flow\\
     Termination mechanic\\
     Gymnasium API};

  \node[stepbox, draw=gray, dashed, right=0.6cm of s2] (s3)
    {\textbf{Step 3}\\[4pt]
     Validated Deployment\\[6pt]
     \scriptsize
     RL agent training\\
     OR benchmark\\
     Primavera~P6 integration\\[4pt]
     \emph{(future work)}};

  \draw[arrow]                    (s1) -- (s2);
  \draw[arrow, dashed, color=gray] (s2) -- (s3);

\end{tikzpicture}
\begin{tikzpicture}[
    node distance  = 0cm,
    stepbox/.style = {rectangle, draw, rounded corners=3pt,
                      minimum width=3.6cm, minimum height=2.8cm,
                      align=center, font=\small, text width=3.2cm},
    arrow/.style   = {->, >=stealth, thick},
  ]

  \node[stepbox, fill=gray!15] (s1) at (0,0)
    {\textbf{Step 1}\\[4pt]
     Domain Grounding\\[6pt]
     \scriptsize
     Contract mechanics\\
     EVM signals\\
     S-curve profiles\\
     Portfolio selection};

  \node[stepbox, fill=gray!8, right=0.6cm of s1] (s2)
    {\textbf{Step 2}\\[1pt]
     \emph{(this work)}\\[4pt]
     Problem Formulation\\[6pt]
     \scriptsize
     MDP environment\\
     FIDIC cash-flow\\
     Termination mechanic\\
     Gymnasium API};

  \node[stepbox, draw=gray, dashed, right=0.6cm of s2] (s3)
    {\textbf{Step 3}\\[4pt]
     Validated Deployment\\[6pt]
     \scriptsize
     RL agent training\\
     OR benchmark\\
     Primavera~P6 integration\\[4pt]
     \emph{(future work)}};

  \draw[arrow]                    (s1) -- (s2);
  \draw[arrow, dashed, color=gray] (s2) -- (s3);

\end{tikzpicture}
\begin{tikzpicture}[
    node distance  = 0cm,
    stepbox/.style = {rectangle, draw, rounded corners=3pt,
                      minimum width=3.6cm, minimum height=2.8cm,
                      align=center, font=\small, text width=3.2cm},
    arrow/.style   = {->, >=stealth, thick},
  ]

  \node[stepbox, fill=gray!15] (s1) at (0,0)
    {\textbf{Step 1}\\[4pt]
     Domain Grounding\\[6pt]
     \scriptsize
     Contract mechanics\\
     EVM signals\\
     S-curve profiles\\
     Portfolio selection};

  \node[stepbox, fill=gray!8, right=0.6cm of s1] (s2)
    {\textbf{Step 2}\\[1pt]
     \emph{(this work)}\\[4pt]
     Problem Formulation\\[6pt]
     \scriptsize
     MDP environment\\
     FIDIC cash-flow\\
     Termination mechanic\\
     Gymnasium API};

  \node[stepbox, draw=gray, dashed, right=0.6cm of s2] (s3)
    {\textbf{Step 3}\\[4pt]
     Validated Deployment\\[6pt]
     \scriptsize
     RL agent training\\
     OR benchmark\\
     Primavera~P6 integration\\[4pt]
     \emph{(future work)}};

  \draw[arrow]                    (s1) -- (s2);
  \draw[arrow, dashed, color=gray] (s2) -- (s3);

\end{tikzpicture}
    \caption{پیشرفت بلوغ به سمت یک سیستم پشتیبانی تصمیم سطح سبد
      قابل استقرار.
      گام ۲، که توسط پژوهش حاضر تکمیل شده، گام ۳ را ممکن می\zwnj{}سازد.
      حاشیه و فلش نقطه\zwnj{}چین کار آینده را نشان می\zwnj{}دهند.}
    \label{fig:maturity}
\end{figure}

\medskip
پژوهش حاضر هر سه شکاف را پوشش می\zwnj{}دهد و گام ۲ این پیشرفت را تکمیل می\zwnj{}کند.
محیط \lr{MDP} سازگار با \lr{Gymnasium}---
با مکانیزم\zwnj{}های قراردادی پارامتری\zwnj{}پذیر،
بازده تصادفی بودجه\zwnj{}به\zwnj{}پیشرفت به\zwnj{}عنوان منبع واحد عدم\zwnj{}قطعیت،
و یک مکانیزم فسخ دوشرطی درون\zwnj{}زا---
سهم روش\zwnj{}شناختی اصلی را تشکیل می\zwnj{}دهد.
اکتشافی\zwnj{}های قطعی حریصانه در برابر یک بهینه‌ی اطلاعات\zwnj{}کامل
در بخش~\ref{sec:evaluation} ارزیابی می\zwnj{}شوند
و یک شکاف بهینگی کمّی\zwnj{}شده ارائه می\zwnj{}دهند که به\zwnj{}صورت یکنواخت
با پیچیدگی سبد رشد می\zwnj{}کند و هدف بهبود را
برای عامل\zwnj{}های \lr{RL} آینده تعریف می\zwnj{}کند.

\begin{table}[htbp]
    \centering
    \caption{آثار نماینده نگاشت\zwnj{}شده در برابر ابعاد مدل\zwnj{}سازی
      چارچوب حاضر.
      \textit{نوع تصمیم}:
      الف\,=\,انتخاب، ز\,=\,زمان\zwnj{}بندی، ب\,=\,بودجه\zwnj{}بندی.
      \checkmark~=~به\zwnj{}طور کامل پوشش\zwnj{}داده\zwnj{}شده؛\;
      $\circ$~=~تا حدی پوشش\zwnj{}داده\zwnj{}شده؛\;
      {--}~=~پوشش\zwnj{}داده\zwnj{}نشده.}
    \label{tab:litmap}
    % tables/litmap.tex
% tabular only — no table environment, no caption, no label.
% Include in section file as:
%   % tables/litmap.tex
% tabular only — no table environment, no caption, no label.
% Include in section file as:
%   % tables/litmap.tex
% tabular only — no table environment, no caption, no label.
% Include in section file as:
%   \input{tables/litmap.tex}
\setlength{\tabcolsep}{6pt}
\renewcommand{\arraystretch}{1.35}
\rowcolors{3}{rowodd}{white}
\resizebox{\textwidth}{!}{%
\begin{tabular}{>{\raggedright\arraybackslash}p{3.8cm}
                c
                c
                >{\columncolor{rarebg}}c
                c
                >{\columncolor{rarebg}}c
                >{\columncolor{rarebg}}c}
  \toprule
  \rowcolor{headbg}
    \color{headfg}\textbf{Work}
    & \color{headfg}\textbf{Portfolio}
    & \color{headfg}\textbf{Decision}
    & \cellcolor{headbg}\color{headfg}\textbf{Sequential}
    & \color{headfg}\textbf{Stochastic}
    & \cellcolor{headbg}\color{headfg}\textbf{Contract}
    & \cellcolor{headbg}\color{headfg}\textbf{Termination} \\
  \rowcolor{headbg}
    \color{headfg}\textbf{}
    & \color{headfg}\textbf{level}
    & \color{headfg}\textbf{type}
    & \cellcolor{headbg}\color{headfg}\textbf{decisions}
    & \color{headfg}\textbf{cash flows}
    & \cellcolor{headbg}\color{headfg}\textbf{mechanics}
    & \cellcolor{headbg}\color{headfg}\textbf{settlement} \\
  \midrule
  \citet{Yamakawa2019}        & \textcolor{symgreen}{\checkmark} & \textsc{\small S}    & \textcolor{symgrey}{--}          & \textcolor{symamber}{$\circ$}    & \textcolor{symgrey}{--}          & \textcolor{symgrey}{--}          \\
  \citet{KOSZTYAN2022}        & \textcolor{symgreen}{\checkmark} & \textsc{\small S}    & \textcolor{symgrey}{--}          & \textcolor{symamber}{$\circ$}    & \textcolor{symgrey}{--}          & \textcolor{symgrey}{--}          \\
  \citet{touran2006mathematical}  & \textcolor{symgreen}{\checkmark} & \textsc{\small Sc/B} & \textcolor{symgrey}{--}          & \textcolor{symamber}{$\circ$}    & \textcolor{symgrey}{--}          & \textcolor{symgrey}{--}          \\
  \citet{touran2010probabilistic} & \textcolor{symgreen}{\checkmark} & \textsc{\small Sc/B} & \textcolor{symgrey}{--}          & \textcolor{symamber}{$\circ$}    & \textcolor{symgrey}{--}          & \textcolor{symgrey}{--}          \\
  \citet{GHAPANCHI2012791}    & \textcolor{symgreen}{\checkmark} & \textsc{\small S}    & \textcolor{symgrey}{--}          & \textcolor{symamber}{$\circ$}    & \textcolor{symgrey}{--}          & \textcolor{symgrey}{--}          \\
  \citet{pourahmadi2015scenario}  & \textcolor{symgreen}{\checkmark} & \textsc{\small S}    & \textcolor{symgrey}{--}          & \textcolor{symamber}{$\circ$}    & \textcolor{symgrey}{--}          & \textcolor{symgrey}{--}          \\
  \citet{elsabek2018coordination} & \textcolor{symgrey}{--}          & \textsc{\small B}    & \textcolor{symgrey}{--}          & \textcolor{symgrey}{--}          & \textcolor{symgreen}{\checkmark} & \textcolor{symgrey}{--}          \\
  \citet{beseiso2021}         & \textcolor{symgreen}{\checkmark} & \textsc{\small S}    & \textcolor{symgrey}{--}          & \textcolor{symamber}{$\circ$}    & \textcolor{symgrey}{--}          & \textcolor{symgrey}{--}          \\
  \citet{karimi2021scenario}  & \textcolor{symgreen}{\checkmark} & \textsc{\small S}    & \textcolor{symgrey}{--}          & \textcolor{symamber}{$\circ$}    & \textcolor{symgrey}{--}          & \textcolor{symgrey}{--}          \\
  \citet{bahroun2023artificial}   & \textcolor{symgrey}{--}          & \textsc{\small Sc}   & \textcolor{symgreen}{\checkmark} & \textcolor{symgrey}{--}          & \textcolor{symgrey}{--}          & \textcolor{symgrey}{--}          \\
  \citet{RANJBAR2022116134}   & \textcolor{symgreen}{\checkmark} & \textsc{\small S/Sc} & \textcolor{symgrey}{--}          & \textcolor{symgrey}{--}          & \textcolor{symamber}{$\circ$}    & \textcolor{symgrey}{--}          \\
  \citet{szwarcfiter2023balancing}& \textcolor{symgreen}{\checkmark} & \textsc{\small Sc}   & \textcolor{symgreen}{\checkmark} & \textcolor{symamber}{$\circ$}    & \textcolor{symgrey}{--}          & \textcolor{symgrey}{--}          \\
  \citet{asghari2023}         & \textcolor{symgrey}{--}          & \textcolor{symgrey}{--} & \textcolor{symgreen}{\checkmark} & \textcolor{symamber}{$\circ$} & \textcolor{symgrey}{--}          & \textcolor{symgrey}{--}          \\
  \citet{mazhari2024}         & \textcolor{symgreen}{\checkmark} & \textsc{\small S/B}  & \textcolor{symgreen}{\checkmark} & \textcolor{symgrey}{--}          & \textcolor{symgrey}{--}          & \textcolor{symgrey}{--}          \\
  \citet{JIANG2024101016}     & \textcolor{symgreen}{\checkmark} & \textcolor{symgrey}{--} & \textcolor{symgreen}{\checkmark} & \textcolor{symamber}{$\circ$} & \textcolor{symgrey}{--}          & \textcolor{symgrey}{--}          \\
  \citet{Nabot2025}           & \textcolor{symgreen}{\checkmark} & \textsc{\small Sc}   & \textcolor{symgreen}{\checkmark} & \textcolor{symgrey}{--}          & \textcolor{symgreen}{\checkmark} & \textcolor{symgreen}{\checkmark} \\
  \citet{Darvish2025}         & \textcolor{symgreen}{\checkmark} & \textsc{\small B}    & \textcolor{symgreen}{\checkmark} & \textcolor{symamber}{$\circ$}    & \textcolor{symamber}{$\circ$}    & \textcolor{symgreen}{\checkmark} \\
  \citet{Gholami2025}         & \textcolor{symgreen}{\checkmark} & \textsc{\small Sc}   & \textcolor{symgreen}{\checkmark} & \textcolor{symamber}{$\circ$}    & \textcolor{symgrey}{--}          & \textcolor{symgreen}{\checkmark} \\
  \citet{Kehinde2025}         & \textcolor{symgreen}{\checkmark} & \textsc{\small B}    & \textcolor{symgreen}{\checkmark} & \textcolor{symgreen}{\checkmark} & \textcolor{symgrey}{--}          & \textcolor{symgrey}{--}          \\
  \citet{Ajibade2025}         & \textcolor{symgreen}{\checkmark} & \textsc{\small Sc}   & \textcolor{symgreen}{\checkmark} & \textcolor{symamber}{$\circ$}    & \textcolor{symgrey}{--}          & \textcolor{symgrey}{--}          \\
  \citet{Huddar2025}          & \textcolor{symgreen}{\checkmark} & \textsc{\small B}    & \textcolor{symgreen}{\checkmark} & \textcolor{symamber}{$\circ$}    & \textcolor{symgrey}{--}          & \textcolor{symgrey}{--}          \\
  \citet{Akbari2026}          & \textcolor{symgreen}{\checkmark} & \textsc{\small B}    & \textcolor{symgreen}{\checkmark} & \textcolor{symamber}{$\circ$}    & \textcolor{symgreen}{\checkmark} & \textcolor{symgreen}{\checkmark} \\
  \citet{ebnerasoul2023}      & \textcolor{symgreen}{\checkmark} & \textsc{\small S}    & \textcolor{symgrey}{--}          & \textcolor{symamber}{$\circ$}    & \textcolor{symgrey}{--}          & \textcolor{symgrey}{--}          \\
  \citet{noroozi2016probabilistic}& \textcolor{symgreen}{\checkmark} & \textsc{\small B}  & \textcolor{symgrey}{--}          & \textcolor{symamber}{$\circ$}    & \textcolor{symgrey}{--}          & \textcolor{symgrey}{--}          \\
  \citet{JIANG2023104817}     & \textcolor{symgrey}{--}          & \textsc{\small B}    & \textcolor{symgreen}{\checkmark} & \textcolor{symgrey}{--}          & \textcolor{symgrey}{--}          & \textcolor{symgrey}{--}          \\
  \cmidrule{1-7}
  \rowcolor{pwbg}
    \color{headfg}\textbf{Present work}
    & \cellcolor{pwbg}\color{headfg}\checkmark
    & \color{headfg}\textbf{B}
    & \cellcolor{pwbg}\color{headfg}\checkmark
    & \color{headfg}\checkmark
    & \cellcolor{pwbg}\color{headfg}\checkmark
    & \cellcolor{pwbg}\color{headfg}\checkmark \\
  \bottomrule
\end{tabular}%
}
\setlength{\tabcolsep}{6pt}
\renewcommand{\arraystretch}{1.35}
\rowcolors{3}{rowodd}{white}
\resizebox{\textwidth}{!}{%
\begin{tabular}{>{\raggedright\arraybackslash}p{3.8cm}
                c
                c
                >{\columncolor{rarebg}}c
                c
                >{\columncolor{rarebg}}c
                >{\columncolor{rarebg}}c}
  \toprule
  \rowcolor{headbg}
    \color{headfg}\textbf{Work}
    & \color{headfg}\textbf{Portfolio}
    & \color{headfg}\textbf{Decision}
    & \cellcolor{headbg}\color{headfg}\textbf{Sequential}
    & \color{headfg}\textbf{Stochastic}
    & \cellcolor{headbg}\color{headfg}\textbf{Contract}
    & \cellcolor{headbg}\color{headfg}\textbf{Termination} \\
  \rowcolor{headbg}
    \color{headfg}\textbf{}
    & \color{headfg}\textbf{level}
    & \color{headfg}\textbf{type}
    & \cellcolor{headbg}\color{headfg}\textbf{decisions}
    & \color{headfg}\textbf{cash flows}
    & \cellcolor{headbg}\color{headfg}\textbf{mechanics}
    & \cellcolor{headbg}\color{headfg}\textbf{settlement} \\
  \midrule
  \citet{Yamakawa2019}        & \textcolor{symgreen}{\checkmark} & \textsc{\small S}    & \textcolor{symgrey}{--}          & \textcolor{symamber}{$\circ$}    & \textcolor{symgrey}{--}          & \textcolor{symgrey}{--}          \\
  \citet{KOSZTYAN2022}        & \textcolor{symgreen}{\checkmark} & \textsc{\small S}    & \textcolor{symgrey}{--}          & \textcolor{symamber}{$\circ$}    & \textcolor{symgrey}{--}          & \textcolor{symgrey}{--}          \\
  \citet{touran2006mathematical}  & \textcolor{symgreen}{\checkmark} & \textsc{\small Sc/B} & \textcolor{symgrey}{--}          & \textcolor{symamber}{$\circ$}    & \textcolor{symgrey}{--}          & \textcolor{symgrey}{--}          \\
  \citet{touran2010probabilistic} & \textcolor{symgreen}{\checkmark} & \textsc{\small Sc/B} & \textcolor{symgrey}{--}          & \textcolor{symamber}{$\circ$}    & \textcolor{symgrey}{--}          & \textcolor{symgrey}{--}          \\
  \citet{GHAPANCHI2012791}    & \textcolor{symgreen}{\checkmark} & \textsc{\small S}    & \textcolor{symgrey}{--}          & \textcolor{symamber}{$\circ$}    & \textcolor{symgrey}{--}          & \textcolor{symgrey}{--}          \\
  \citet{pourahmadi2015scenario}  & \textcolor{symgreen}{\checkmark} & \textsc{\small S}    & \textcolor{symgrey}{--}          & \textcolor{symamber}{$\circ$}    & \textcolor{symgrey}{--}          & \textcolor{symgrey}{--}          \\
  \citet{elsabek2018coordination} & \textcolor{symgrey}{--}          & \textsc{\small B}    & \textcolor{symgrey}{--}          & \textcolor{symgrey}{--}          & \textcolor{symgreen}{\checkmark} & \textcolor{symgrey}{--}          \\
  \citet{beseiso2021}         & \textcolor{symgreen}{\checkmark} & \textsc{\small S}    & \textcolor{symgrey}{--}          & \textcolor{symamber}{$\circ$}    & \textcolor{symgrey}{--}          & \textcolor{symgrey}{--}          \\
  \citet{karimi2021scenario}  & \textcolor{symgreen}{\checkmark} & \textsc{\small S}    & \textcolor{symgrey}{--}          & \textcolor{symamber}{$\circ$}    & \textcolor{symgrey}{--}          & \textcolor{symgrey}{--}          \\
  \citet{bahroun2023artificial}   & \textcolor{symgrey}{--}          & \textsc{\small Sc}   & \textcolor{symgreen}{\checkmark} & \textcolor{symgrey}{--}          & \textcolor{symgrey}{--}          & \textcolor{symgrey}{--}          \\
  \citet{RANJBAR2022116134}   & \textcolor{symgreen}{\checkmark} & \textsc{\small S/Sc} & \textcolor{symgrey}{--}          & \textcolor{symgrey}{--}          & \textcolor{symamber}{$\circ$}    & \textcolor{symgrey}{--}          \\
  \citet{szwarcfiter2023balancing}& \textcolor{symgreen}{\checkmark} & \textsc{\small Sc}   & \textcolor{symgreen}{\checkmark} & \textcolor{symamber}{$\circ$}    & \textcolor{symgrey}{--}          & \textcolor{symgrey}{--}          \\
  \citet{asghari2023}         & \textcolor{symgrey}{--}          & \textcolor{symgrey}{--} & \textcolor{symgreen}{\checkmark} & \textcolor{symamber}{$\circ$} & \textcolor{symgrey}{--}          & \textcolor{symgrey}{--}          \\
  \citet{mazhari2024}         & \textcolor{symgreen}{\checkmark} & \textsc{\small S/B}  & \textcolor{symgreen}{\checkmark} & \textcolor{symgrey}{--}          & \textcolor{symgrey}{--}          & \textcolor{symgrey}{--}          \\
  \citet{JIANG2024101016}     & \textcolor{symgreen}{\checkmark} & \textcolor{symgrey}{--} & \textcolor{symgreen}{\checkmark} & \textcolor{symamber}{$\circ$} & \textcolor{symgrey}{--}          & \textcolor{symgrey}{--}          \\
  \citet{Nabot2025}           & \textcolor{symgreen}{\checkmark} & \textsc{\small Sc}   & \textcolor{symgreen}{\checkmark} & \textcolor{symgrey}{--}          & \textcolor{symgreen}{\checkmark} & \textcolor{symgreen}{\checkmark} \\
  \citet{Darvish2025}         & \textcolor{symgreen}{\checkmark} & \textsc{\small B}    & \textcolor{symgreen}{\checkmark} & \textcolor{symamber}{$\circ$}    & \textcolor{symamber}{$\circ$}    & \textcolor{symgreen}{\checkmark} \\
  \citet{Gholami2025}         & \textcolor{symgreen}{\checkmark} & \textsc{\small Sc}   & \textcolor{symgreen}{\checkmark} & \textcolor{symamber}{$\circ$}    & \textcolor{symgrey}{--}          & \textcolor{symgreen}{\checkmark} \\
  \citet{Kehinde2025}         & \textcolor{symgreen}{\checkmark} & \textsc{\small B}    & \textcolor{symgreen}{\checkmark} & \textcolor{symgreen}{\checkmark} & \textcolor{symgrey}{--}          & \textcolor{symgrey}{--}          \\
  \citet{Ajibade2025}         & \textcolor{symgreen}{\checkmark} & \textsc{\small Sc}   & \textcolor{symgreen}{\checkmark} & \textcolor{symamber}{$\circ$}    & \textcolor{symgrey}{--}          & \textcolor{symgrey}{--}          \\
  \citet{Huddar2025}          & \textcolor{symgreen}{\checkmark} & \textsc{\small B}    & \textcolor{symgreen}{\checkmark} & \textcolor{symamber}{$\circ$}    & \textcolor{symgrey}{--}          & \textcolor{symgrey}{--}          \\
  \citet{Akbari2026}          & \textcolor{symgreen}{\checkmark} & \textsc{\small B}    & \textcolor{symgreen}{\checkmark} & \textcolor{symamber}{$\circ$}    & \textcolor{symgreen}{\checkmark} & \textcolor{symgreen}{\checkmark} \\
  \citet{ebnerasoul2023}      & \textcolor{symgreen}{\checkmark} & \textsc{\small S}    & \textcolor{symgrey}{--}          & \textcolor{symamber}{$\circ$}    & \textcolor{symgrey}{--}          & \textcolor{symgrey}{--}          \\
  \citet{noroozi2016probabilistic}& \textcolor{symgreen}{\checkmark} & \textsc{\small B}  & \textcolor{symgrey}{--}          & \textcolor{symamber}{$\circ$}    & \textcolor{symgrey}{--}          & \textcolor{symgrey}{--}          \\
  \citet{JIANG2023104817}     & \textcolor{symgrey}{--}          & \textsc{\small B}    & \textcolor{symgreen}{\checkmark} & \textcolor{symgrey}{--}          & \textcolor{symgrey}{--}          & \textcolor{symgrey}{--}          \\
  \cmidrule{1-7}
  \rowcolor{pwbg}
    \color{headfg}\textbf{Present work}
    & \cellcolor{pwbg}\color{headfg}\checkmark
    & \color{headfg}\textbf{B}
    & \cellcolor{pwbg}\color{headfg}\checkmark
    & \color{headfg}\checkmark
    & \cellcolor{pwbg}\color{headfg}\checkmark
    & \cellcolor{pwbg}\color{headfg}\checkmark \\
  \bottomrule
\end{tabular}%
}
\setlength{\tabcolsep}{6pt}
\renewcommand{\arraystretch}{1.35}
\rowcolors{3}{rowodd}{white}
\resizebox{\textwidth}{!}{%
\begin{tabular}{>{\raggedright\arraybackslash}p{3.8cm}
                c
                c
                >{\columncolor{rarebg}}c
                c
                >{\columncolor{rarebg}}c
                >{\columncolor{rarebg}}c}
  \toprule
  \rowcolor{headbg}
    \color{headfg}\textbf{Work}
    & \color{headfg}\textbf{Portfolio}
    & \color{headfg}\textbf{Decision}
    & \cellcolor{headbg}\color{headfg}\textbf{Sequential}
    & \color{headfg}\textbf{Stochastic}
    & \cellcolor{headbg}\color{headfg}\textbf{Contract}
    & \cellcolor{headbg}\color{headfg}\textbf{Termination} \\
  \rowcolor{headbg}
    \color{headfg}\textbf{}
    & \color{headfg}\textbf{level}
    & \color{headfg}\textbf{type}
    & \cellcolor{headbg}\color{headfg}\textbf{decisions}
    & \color{headfg}\textbf{cash flows}
    & \cellcolor{headbg}\color{headfg}\textbf{mechanics}
    & \cellcolor{headbg}\color{headfg}\textbf{settlement} \\
  \midrule
  \citet{Yamakawa2019}        & \textcolor{symgreen}{\checkmark} & \textsc{\small S}    & \textcolor{symgrey}{--}          & \textcolor{symamber}{$\circ$}    & \textcolor{symgrey}{--}          & \textcolor{symgrey}{--}          \\
  \citet{KOSZTYAN2022}        & \textcolor{symgreen}{\checkmark} & \textsc{\small S}    & \textcolor{symgrey}{--}          & \textcolor{symamber}{$\circ$}    & \textcolor{symgrey}{--}          & \textcolor{symgrey}{--}          \\
  \citet{touran2006mathematical}  & \textcolor{symgreen}{\checkmark} & \textsc{\small Sc/B} & \textcolor{symgrey}{--}          & \textcolor{symamber}{$\circ$}    & \textcolor{symgrey}{--}          & \textcolor{symgrey}{--}          \\
  \citet{touran2010probabilistic} & \textcolor{symgreen}{\checkmark} & \textsc{\small Sc/B} & \textcolor{symgrey}{--}          & \textcolor{symamber}{$\circ$}    & \textcolor{symgrey}{--}          & \textcolor{symgrey}{--}          \\
  \citet{GHAPANCHI2012791}    & \textcolor{symgreen}{\checkmark} & \textsc{\small S}    & \textcolor{symgrey}{--}          & \textcolor{symamber}{$\circ$}    & \textcolor{symgrey}{--}          & \textcolor{symgrey}{--}          \\
  \citet{pourahmadi2015scenario}  & \textcolor{symgreen}{\checkmark} & \textsc{\small S}    & \textcolor{symgrey}{--}          & \textcolor{symamber}{$\circ$}    & \textcolor{symgrey}{--}          & \textcolor{symgrey}{--}          \\
  \citet{elsabek2018coordination} & \textcolor{symgrey}{--}          & \textsc{\small B}    & \textcolor{symgrey}{--}          & \textcolor{symgrey}{--}          & \textcolor{symgreen}{\checkmark} & \textcolor{symgrey}{--}          \\
  \citet{beseiso2021}         & \textcolor{symgreen}{\checkmark} & \textsc{\small S}    & \textcolor{symgrey}{--}          & \textcolor{symamber}{$\circ$}    & \textcolor{symgrey}{--}          & \textcolor{symgrey}{--}          \\
  \citet{karimi2021scenario}  & \textcolor{symgreen}{\checkmark} & \textsc{\small S}    & \textcolor{symgrey}{--}          & \textcolor{symamber}{$\circ$}    & \textcolor{symgrey}{--}          & \textcolor{symgrey}{--}          \\
  \citet{bahroun2023artificial}   & \textcolor{symgrey}{--}          & \textsc{\small Sc}   & \textcolor{symgreen}{\checkmark} & \textcolor{symgrey}{--}          & \textcolor{symgrey}{--}          & \textcolor{symgrey}{--}          \\
  \citet{RANJBAR2022116134}   & \textcolor{symgreen}{\checkmark} & \textsc{\small S/Sc} & \textcolor{symgrey}{--}          & \textcolor{symgrey}{--}          & \textcolor{symamber}{$\circ$}    & \textcolor{symgrey}{--}          \\
  \citet{szwarcfiter2023balancing}& \textcolor{symgreen}{\checkmark} & \textsc{\small Sc}   & \textcolor{symgreen}{\checkmark} & \textcolor{symamber}{$\circ$}    & \textcolor{symgrey}{--}          & \textcolor{symgrey}{--}          \\
  \citet{asghari2023}         & \textcolor{symgrey}{--}          & \textcolor{symgrey}{--} & \textcolor{symgreen}{\checkmark} & \textcolor{symamber}{$\circ$} & \textcolor{symgrey}{--}          & \textcolor{symgrey}{--}          \\
  \citet{mazhari2024}         & \textcolor{symgreen}{\checkmark} & \textsc{\small S/B}  & \textcolor{symgreen}{\checkmark} & \textcolor{symgrey}{--}          & \textcolor{symgrey}{--}          & \textcolor{symgrey}{--}          \\
  \citet{JIANG2024101016}     & \textcolor{symgreen}{\checkmark} & \textcolor{symgrey}{--} & \textcolor{symgreen}{\checkmark} & \textcolor{symamber}{$\circ$} & \textcolor{symgrey}{--}          & \textcolor{symgrey}{--}          \\
  \citet{Nabot2025}           & \textcolor{symgreen}{\checkmark} & \textsc{\small Sc}   & \textcolor{symgreen}{\checkmark} & \textcolor{symgrey}{--}          & \textcolor{symgreen}{\checkmark} & \textcolor{symgreen}{\checkmark} \\
  \citet{Darvish2025}         & \textcolor{symgreen}{\checkmark} & \textsc{\small B}    & \textcolor{symgreen}{\checkmark} & \textcolor{symamber}{$\circ$}    & \textcolor{symamber}{$\circ$}    & \textcolor{symgreen}{\checkmark} \\
  \citet{Gholami2025}         & \textcolor{symgreen}{\checkmark} & \textsc{\small Sc}   & \textcolor{symgreen}{\checkmark} & \textcolor{symamber}{$\circ$}    & \textcolor{symgrey}{--}          & \textcolor{symgreen}{\checkmark} \\
  \citet{Kehinde2025}         & \textcolor{symgreen}{\checkmark} & \textsc{\small B}    & \textcolor{symgreen}{\checkmark} & \textcolor{symgreen}{\checkmark} & \textcolor{symgrey}{--}          & \textcolor{symgrey}{--}          \\
  \citet{Ajibade2025}         & \textcolor{symgreen}{\checkmark} & \textsc{\small Sc}   & \textcolor{symgreen}{\checkmark} & \textcolor{symamber}{$\circ$}    & \textcolor{symgrey}{--}          & \textcolor{symgrey}{--}          \\
  \citet{Huddar2025}          & \textcolor{symgreen}{\checkmark} & \textsc{\small B}    & \textcolor{symgreen}{\checkmark} & \textcolor{symamber}{$\circ$}    & \textcolor{symgrey}{--}          & \textcolor{symgrey}{--}          \\
  \citet{Akbari2026}          & \textcolor{symgreen}{\checkmark} & \textsc{\small B}    & \textcolor{symgreen}{\checkmark} & \textcolor{symamber}{$\circ$}    & \textcolor{symgreen}{\checkmark} & \textcolor{symgreen}{\checkmark} \\
  \citet{ebnerasoul2023}      & \textcolor{symgreen}{\checkmark} & \textsc{\small S}    & \textcolor{symgrey}{--}          & \textcolor{symamber}{$\circ$}    & \textcolor{symgrey}{--}          & \textcolor{symgrey}{--}          \\
  \citet{noroozi2016probabilistic}& \textcolor{symgreen}{\checkmark} & \textsc{\small B}  & \textcolor{symgrey}{--}          & \textcolor{symamber}{$\circ$}    & \textcolor{symgrey}{--}          & \textcolor{symgrey}{--}          \\
  \citet{JIANG2023104817}     & \textcolor{symgrey}{--}          & \textsc{\small B}    & \textcolor{symgreen}{\checkmark} & \textcolor{symgrey}{--}          & \textcolor{symgrey}{--}          & \textcolor{symgrey}{--}          \\
  \cmidrule{1-7}
  \rowcolor{pwbg}
    \color{headfg}\textbf{Present work}
    & \cellcolor{pwbg}\color{headfg}\checkmark
    & \color{headfg}\textbf{B}
    & \cellcolor{pwbg}\color{headfg}\checkmark
    & \color{headfg}\checkmark
    & \cellcolor{pwbg}\color{headfg}\checkmark
    & \cellcolor{pwbg}\color{headfg}\checkmark \\
  \bottomrule
\end{tabular}%
}
\end{table}